% !TEX TS-program = xelatex
% !TEX encoding = UTF-8 Unicode
% !Mode:: "TeX:UTF-8"

\documentclass{resume}
\usepackage{zh_CN-Adobefonts_external} % Simplified Chinese Support using external fonts (./fonts/zh_CN-Adobe/)
%\usepackage{zh_CN-Adobefonts_internal} % Simplified Chinese Support using system fonts
\usepackage{linespacing_fix} % disable extra space before next section
\usepackage{cite}
\usepackage{xcolor}

\begin{document}
\pagenumbering{gobble} % suppress displaying page number


\name{姜海涛}
\vspace{4pt}

% {
%   \centering
%   \LARGE\bfseries 姜海涛\par
% }
% \vspace{16pt} % 或你想要的任意间距


\basicInfo{
  \email{2125996426@qq.com} \textperiodcentered\ 
  \phone{(+86) 135-2376-8662} \textperiodcentered\ 
  \faBullseye{\small 求职意向:算法工程师}
}


\vspace{-0.6em}
% ######## education #########
\section{\faGraduationCap\  教育背景}
\vspace{-0.2em}
\datedsubsection{\textbf{西安交通大学}\quad \textit{人工智能}\quad \textit{在读硕士}}{2024 - Present}
% \begin{itemize}
%   \item 研究方向：视觉语言模型（VLM）、大语言模型（LLM）
% \end{itemize}

% \vspace{-0.8em}
\datedsubsection{\textbf{杭州电子科技大学}\quad \textit{机械设计制造及其自动化}\quad \textit{学士}}{2019 - 2023}
% \begin{itemize}
%   \item 相关课程：数理统计、数据结构与算法、操作系统、计算机网络
% \end{itemize}


% ######## internship ########
\section{\faLaptop\ 实习经历}
\vspace{-0.2em}

% % ######## 生成-评估 重排 ########
% \datedsubsection{\textbf{快手科技}\quad \textit{主站推荐中心 · 算法工程师}}{2026.06 - Present}
% \role{负责快手主站双列推荐 Generator--Evaluator 重排链路优化，研究列表级价值建模与生成评估协同优化，提升重排系统对高价值消费序列的识别和搜索能力}{}
% \vspace{-0.6em}
% \begin{itemize}
%   \item \textbf{Evaluator 优化:}
%   将现有逐 item 预估、规则累加的评估范式升级为消费过程感知的 List-wise Model。基于 Causal Prefix Transformer 建模序列上下文，将消费长度分解为逐位置继续概率并递推 $P(K)$，联合预测各 Prefix 的累计观看时长、有效 VV 与互动价值；针对曝光截断、长尾价值标签及未曝光后缀缺少反事实监督等问题，设计右截断 Prefix / List 双层价值回归、累计价值单调约束与反馈偏好 RankNet Loss，最终通过
%   $\sum_k P(K=k)V(\mathrm{Prefix}_{1:k})$
%   直接估计候选 List 的长期消费价值。

%   \item \textbf{G--E 联合优化:}
%   完成 Evaluator-guided Generator 自迭代方案设计。借鉴 NLGR 的邻域相对奖励与 DeGRe 的 Lookahead Search，由 Evaluator 离线挖掘高价值 Beam Paths，并向 AR/NAR Generator 蒸馏列表偏好、路径权重与 Step-wise Action Distribution，推进“候选生成—列表评估—稠密监督—生成器迭代”的优化闭环。
% \end{itemize}


% ######## 生成式召回 ########
% \datedsubsection{\textbf{快手科技}\quad \textit{主站推荐中心 · 推荐大模型算法工程师}}{2026.06 - Present}
% \role{负责快手主站双列外流生成式召回通道建设，针对 SID 自回归生成中的语义漂移与线上时延问题，完成轻量生成架构、路径级 PRM 及推理链路优化，实现增量候选召回与线上消费提升}{}
% \vspace{-0.6em}
% \begin{itemize}
%   \item \textbf{轻量生成式召回:}
%   基于用户画像与最长 200 条点击序列，自回归生成三层 Semantic ID，并通过 SID 索引召回真实视频；将原有独立 4 层 Encoder--Decoder 改造为 MLP + LayerNorm 用户表征与轻量 Decoder 架构，由 Decoder Cross-Attention 直接读取用户上下文，降低生成式召回的线上计算开销。

%   \item \textbf{路径级过程监督:}
%   针对 Teacher Forcing 与自回归推理不一致造成的误差累积，以及 NTP 局部 Token 目标难以刻画完整路径用户匹配度的问题，引入路径级 PRM，将各层 SID Prefix 视为中间生成步骤，通过 NTP + InfoNCE 端到端联合训练提供逐层稠密监督；设计 False Negative Mask 排除同路径假负样本，并采用 logQ Correction 缓解热门路径采样偏差。

%   \item \textbf{PRM 引导搜索与线上收益:}
%   设计 PRM-guided Beam Search，由 Decoder 按累积生成概率预筛 Top-1024 路径，再由 PRM 剪枝保留 Top-512，及时阻断低匹配度分支、缓解路径语义漂移；结合高质量消费样本过滤、播放时长加权、多级 KV Cache 与混合精度部署，在 Leaf 平均耗时仅增加 1.411\% 的情况下，召回候选独占率达到 31.43\%，发现页外流播放时长与次均播放时长分别提升 0.358\%/0.409\%，关注率与收藏率分别提升 0.892\%/0.794\%，负反馈下降 1.415\%。
% \end{itemize}

\datedsubsection{\textbf{快手科技}\quad \textit{主站推荐中心 · 推荐大模型算法工程师}}{2026.06 - Present}
\role{负责快手主站双列推荐生成式召回优化，针对 SID 自回归生成中的路径语义漂移与线上时延问题，完成轻量生成架构、路径级 PRM 及推理链路优化，实现增量候选召回与线上消费提升}{}
\vspace{-0.6em}
\begin{itemize}
  \item \textbf{生成式架构:}
  基于用户画像与最长 200 条行为序列，自回归生成三层 Semantic ID，并通过 SID 索引完成真实视频召回；将原有独立 4 层 Encoder--Decoder 改造为轻量 Decoder 架构，由 MLP 编码用户特征并直接作为 Cross-Attention 的 Key / Value，降低序列建模与线上生成开销。

  \item \textbf{路径级 PRM:}
  % 针对 NTP 仅优化局部 Token 概率、难以衡量完整路径与用户兴趣匹配度的问题，设计轻量级路径判别模型；
  将各层 SID Prefix Embedding 累加为路径 Query，通过 Cross-Attention 聚合用户上下文并由 MLP 输出用户--路径匹配分数；以 Batch 内其他用户的同层路径作为负样本，通过 NTP + InfoNCE 端到端联合训练，为每一步 SID 生成提供路径级监督；设计 False Negative Mask 排除重复路径假负样本，并采用 logQ Correction 缓解热门路径采样偏差。

  \item \textbf{推理部署优化:}
  由 Decoder 按累积生成概率预筛 Top-1024 路径，再利用 PRM 匹配分数剪枝至 Top-512，及时淘汰偏离用户兴趣的中间分支；实现 Decoder Self / Cross-Attention 与 PRM 多级 KV Cache，结合 FP16 / FP32 混合精度，控制大规模 Beam Search 的延迟与显存开销。

  \item \textbf{线上收益:}
  召回候选独占率达到 31.43\%，透出占比约 9.914\%；在 Leaf 平均耗时仅增加 1.411\% 的情况下，外流播放时长提升 0.358\%。
\end{itemize}


\vspace{-1.0em}
% ######## 项目经历 ########
\section{\faBriefcase\ 项目经历}
\vspace{-0.2em}

% \faBriefcase \faBuilding \faIndustry \faSuitcase
% \faUsers \faCogs \faCode \faRocket \faFlask \faTasks \faLaptop \faFolderOpen \faLightbulbO

\datedsubsection{\textbf{通用领域大模型推理增强} 
\quad\quad\quad
\textit{概率奖励驱动的无验证器后训练}}{2025.09 - 2025.12}
\begin{itemize}
  % \item \textbf{项目描述:} 
  % 面向复杂场景下逻辑推演的需求，整体采用\textbf{冷启动 SFT + RLPR 两阶段训练范式}。
  % \item \textbf{训练数据:} 
  % 基于 DeepSeek-R1 范式构造 8K CoT 数据完成冷启动，使模型掌握规范推理格式。
  % 基于 WebInstruct 数据集，使用大模型进行推理复杂度评分（1-4分），筛选出 77K 高质量 RL 训练样本（评分$\geq$3），采用 10-gram 去重避免数据污染。
  % \item \textbf{技术实现:} 
  % 采用 Token 平均解码概率的内在奖励机制（Probability as Reward）替代传统验证器，避免序列似然的高方差问题，引入奖励去偏策略消除题目本身难度偏差，结合自适应标准差过滤动态筛选样本，解决奖励稀疏导致的训练不稳定。
  % 基于 Qwen2.5-7B 模型，GRPO 算法（KL\_coef=0, clip\_eps=0.27）进行训练。在 MMLU-Pro（+4.2\%）、TheoremQA（+10.1\%）等多个基准上评估，平均推理能力达到 59.7\%（+7.5\%）。

  \item \textbf{项目描述:}
  面向缺少规则验证器的通用领域推理任务，基于 Qwen2.5-7B 构建冷启动 SFT + RLPR 两阶段后训练框架，重点解决传统 RLVR 依赖领域专用验证器、难以扩展至开放域推理的问题。

  \item \textbf{数据构建:}
  基于 DeepSeek-R1 范式蒸馏 8K 高质量 CoT 数据完成冷启动；从 WebInstruct 中按推理复杂度筛选 77K RL 训练样本，并通过 10-gram 去重控制评测集污染。

  \item \textbf{技术实现:}
  将 rollout 生成答案替换为参考答案，以模型在给定推理过程下对参考答案 token 的平均概率构造连续奖励；引入无 CoT 基线去偏消除题目本身难度偏差，结合自适应标准差过滤动态筛选样本，实现更稳定的 GRPO 训练；在 MMLU-Pro、TheoremQA 等基准上平均推理能力达到 59.7\%，较基座提升 7.5\%。
\end{itemize}


\vspace{-0.6em}
\datedsubsection{\textbf{长程工具智能体强化学习}
\quad\quad\quad
\textit{动作对齐优化与可验证过程奖励}}{2026.05 - 2026.07}
\begin{itemize}
  \item \textbf{项目描述:}
  面向 ALFWorld 长程工具交互任务，基于 Qwen3-4B 构建 Agentic RL 训练框架，重点解决稀疏终局奖励下的动作信用分配与无效探索问题。

  \item \textbf{技术实现:}
  针对 GRPO 轨迹级归因过粗、Token-GAE 的信用传播与裁剪粒度同环境决策错位的问题，采用 CAPO 在动作边界估计 critic、沿交互 step 计算 action-level GAE；基于中心化平方根长度校准构造 action-aware ratio，对完整动作统一加权与裁剪。

  \item \textbf{奖励优化:}
  基于 ALFWorld 合法动作集合与符号状态构建 Environment Verifier，识别非法操作和关键子目标进展，融合 terminal reward、invalid penalty 与 potential-based process reward，降低无效探索并增强长程任务表现；在 Seen / Unseen 任务上取得 93.86\% / 88.67\% 成功率，较 GRPO 基线提升 12.43 / 14.04 个百分点。
\end{itemize}


\vspace{-1.0em}
% ######## paper #########
\section{\faBook\ 科研经历}
\vspace{-0.2em}
% \datedsubsection{\textbf{VR-CLIP: Visual Refinement of CLIP for Zero-Shot Semantic Segmentation}}{}
\role{\textbf{VR-CLIP: Visual Refinement of CLIP for Zero-Shot Semantic Segmentation \quad\quad\quad\quad CVPR 2026 \quad 一作}}{}
\role{现有零样本语义分割过度依赖解码器弥补 CLIP 像素表征不足，我们通过浅层局部空间建模与跨层文本语义校正，在保持视觉--语言对齐能力的前提下兼顾像素判别能力与未见类别泛化。}{}
% \vspace{-0.6em}
% \begin{itemize}
%   \item \textbf{问题背景:} 
%   现有方法普遍依赖更强的解码器来弥补 CLIP 图像级表征的密集判别不足问题，而直接微调视觉编码器的方法又会破坏原始 CLIP 的跨模态对齐能力，导致解码阶段匹配失效。
  
%   \item \textbf{成果描述:} 
%   提出一种 \textbf{单阶段 in-encoder 特征增强框架}。在 CLIP ViT 浅层引入 NSA 强化邻近空间信息，避免破坏原始语义对齐能力，设计 SGR 基于相似度映射将高级语义信息注入视觉 patch 特征，实现残差式层级语义校正，并设计\textbf{跨层参数共享机制}学习共性语义映射，大幅提升对未知类别的泛化能力。
%   在 COCO-Stuff 上取得 hIoU 51.3\%（+2.2\%），同时维持轻量级训练参数（8\%）。
% \end{itemize}


% \vspace{-1.0em}
% \section{\faInfo\ 个人描述}
% % increase linespacing [parsep=0.5ex]
% \vspace{-0.2em}
% 本人心态乐观向上，情绪稳定，具备良好的英文能力（CET-6），掌握 Python / C++ 语言，熟悉 VLM、LLM 后训练领域，具备扎实的数理基础与工程实现能力（数学建模一等奖），良好的团队协作与沟通能力（青协优秀社员），能够快速融入团队并推动工程落地。


\end{document}
